% ==============================================================================
% PRÉAMBULE DU RAPPORT DE STAGE D'INGÉNIEUR (ENSA SAFI / GROUPE OCP)
% Style Académique et Professionnel - 4ème Année Génie Informatique
% ==============================================================================

\usepackage[utf8]{inputenc}
\usepackage[T1]{fontenc}
\usepackage[french]{babel}
\usepackage{lmodern}
\usepackage{microtype}

% --- Géométrie des pages ---
\usepackage[
    a4paper,
    top=2.5cm,
    bottom=2.5cm,
    left=2.8cm,
    right=2.5cm,
    headheight=14pt
]{geometry}

% --- Typographie & Espacement ---
\usepackage{setspace}
\setstretch{1.18}
\usepackage{parskip}
\setlength{\parindent}{0pt}
\setlength{\parskip}{6pt plus 2pt minus 1pt}

% --- Couleurs Professionnelles & Académiques ---
\usepackage[table,xcdraw]{xcolor}
\definecolor{ocpnavy}{RGB}{20, 45, 85}       % Bleu marine institutionnel OCP
\definecolor{ocpblue}{RGB}{30, 90, 150}      % Bleu secondaire
\definecolor{ocpgreen}{RGB}{34, 112, 62}     % Vert OCP
\definecolor{ocporange}{RGB}{210, 105, 30}   % Accent orange / ambre
\definecolor{darkslate}{RGB}{44, 62, 80}     % Texte d'en-têtes
\definecolor{codebg}{RGB}{246, 248, 250}     % Fond des codes
\definecolor{codeframe}{RGB}{218, 222, 229}  % Bordure de code
\definecolor{commentgreen}{RGB}{46, 125, 50} % Commentaires de code
\definecolor{keywordblue}{RGB}{0, 70, 150}   % Mots-clés de code
\definecolor{stringred}{RGB}{163, 21, 21}    % Chaînes de caractères
\definecolor{linegray}{RGB}{120, 130, 140}   % Numéros de ligne

% Couleurs d'états opérationnels
\definecolor{stnominal}{RGB}{46, 125, 50}    % Vert nominal (Running)
\definecolor{stwarning}{RGB}{230, 124, 11}   % Orange warning
\definecolor{stcritical}{RGB}{198, 40, 40}   % Rouge critique
\definecolor{stdegraded}{RGB}{123, 31, 162}  % Violet dégradé
\definecolor{stmuted}{RGB}{100, 116, 139}    % Gris inactif (Stopped)

% --- Tableaux ---
\usepackage{booktabs}
\usepackage{tabularx}
\usepackage{multirow}
\usepackage{makecell}
\usepackage{array}
\newcolumntype{Y}{>{\raggedright\arraybackslash}X}
\newcolumntype{C}{>{\centering\arraybackslash}X}

% --- Graphiques & Flottants ---
\usepackage{graphicx}
\usepackage{caption}
\usepackage{subcaption}
\captionsetup{
    font=small,
    labelfont={bf,color=ocpnavy},
    labelsep=quad,
    justification=centering
}

% --- TikZ & PGFPlots (Diagrammes vectoriels professionnels) ---
\usepackage{tikz}
\usepackage{pgfplots}
\pgfplotsset{compat=1.18}
\usetikzlibrary{
    calc,
    shapes.geometric,
    shapes.misc,
    arrows.meta,
    positioning,
    shadows,
    fit,
    backgrounds,
    matrix,
    chains,
    decorations.pathreplacing
}

% Styles TikZ prédéfinis
\tikzset{
    block/.style={
        rectangle,
        draw=ocpnavy,
        fill=white,
        thick,
        rounded corners=3pt,
        text centered,
        minimum height=2.2em,
        font=\small\sffamily
    },
    blockfill/.style={
        rectangle,
        draw=ocpnavy!80,
        fill=ocpnavy!10,
        thick,
        rounded corners=3pt,
        text centered,
        minimum height=2.2em,
        font=\small\sffamily
    },
    cloudnode/.style={
        ellipse,
        draw=ocpblue,
        fill=ocpblue!10,
        thick,
        text centered,
        font=\small\sffamily
    },
    flowarrow/.style={
        ->,
        >=Stealth,
        thick,
        draw=ocpnavy
    },
    dataarrow/.style={
        ->,
        >=Stealth,
        dashed,
        thick,
        draw=ocpblue
    },
    alertarrow/.style={
        ->,
        >=Stealth,
        very thick,
        draw=stcritical
    }
}

% --- Code Source (Listings) ---
\usepackage{listings}
\lstset{
    backgroundcolor=\color{codebg},
    basicstyle=\ttfamily\footnotesize,
    breaklines=true,
    breakatwhitespace=true,
    captionpos=b,
    commentstyle=\color{commentgreen}\itshape,
    frame=single,
    framerule=0.6pt,
    rulecolor=\color{codeframe},
    keepspaces=true,
    keywordstyle=\color{keywordblue}\bfseries,
    numbers=left,
    numbersep=8pt,
    numberstyle=\tiny\color{linegray},
    showspaces=false,
    showstringspaces=false,
    showtabs=false,
    stepnumber=1,
    stringstyle=\color{stringred},
    tabsize=2,
    xleftmargin=12pt,
    xrightmargin=5pt,
    framexleftmargin=8pt,
    literate=
        {á}{{\'a}}1 {é}{{\'e}}1 {í}{{\'i}}1 {ó}{{\'o}}1 {ú}{{\'u}}1
        {Á}{{\'A}}1 {É}{{\'E}}1 {Í}{{\'I}}1 {Ó}{{\'O}}1 {Ú}{{\'U}}1
        {à}{{\`a}}1 {è}{{\`e}}1 {ì}{{\`i}}1 {ò}{{\`o}}1 {ù}{{\`u}}1
        {À}{{\`A}}1 {È}{{\`E}}1 {Ì}{{\`I}}1 {Ò}{{\`O}}1 {Ù}{{\`U}}1
        {â}{{\^a}}1 {ê}{{\^e}}1 {î}{{\^i}}1 {ô}{{\^o}}1 {û}{{\^u}}1
        {Â}{{\^A}}1 {Ê}{{\^E}}1 {Î}{{\^I}}1 {Ô}{{\^O}}1 {Û}{{\^U}}1
        {ç}{{\c{c}}}1 {Ç}{{\c{C}}}1
}

% Langage C# personnalisé pour listings
\lstdefinelanguage{CSharpCustom}{
    morekeywords={
        abstract, as, base, bool, break, byte, case, catch, char, checked,
        class, const, continue, decimal, default, delegate, do, double, else,
        enum, event, explicit, extern, false, finally, fixed, float, for,
        foreach, goto, if, implicit, in, int, interface, internal, is, lock,
        long, namespace, new, null, object, operator, out, override, params,
        private, protected, public, readonly, record, ref, return, sbyte,
        sealed, short, sizeof, stackalloc, static, string, struct, switch,
        this, throw, true, try, typeof, uint, ulong, unchecked, unsafe,
        ushort, using, virtual, void, volatile, while, async, await, var,
        get, set, init, required, yield
    },
    sensitive=true,
    morecomment=[l]{//},
    morecomment=[s]{/*}{*/},
    morestring=[b]",
    morestring=[d]'
}

% --- En-têtes & Pieds de page (Fancyhdr) ---
\usepackage{fancyhdr}
\pagestyle{fancy}
\fancyhf{}
\renewcommand{\chaptermark}[1]{\markboth{\thechapter.\ #1}{}}
\renewcommand{\sectionmark}[1]{\markright{\thesection\ #1}}
\fancyhead[L]{\small\color{darkslate}\leftmark}
\fancyhead[R]{\small\color{darkslate}\rightmark}
\fancyfoot[C]{\small\thepage}
\renewcommand{\headrulewidth}{0.4pt}
\renewcommand{\footrulewidth}{0pt}

% Style pour les pages de début de chapitre
\fancypagestyle{plain}{
    \fancyhf{}
    \fancyfoot[C]{\small\thepage}
    \renewcommand{\headrulewidth}{0pt}
    \renewcommand{\footrulewidth}{0pt}
}

% --- Titres de Chapitres & Sections ---
\usepackage{titlesec}
\titleformat{\chapter}[display]
    {\normalfont\sffamily\huge\bfseries\color{ocpnavy}}
    {\large\scshape\color{ocpblue}\chaptertitlename\ \thechapter}
    {6pt}
    {\titlerule[1.2pt]\vspace{3pt}}
    [\vspace{3pt}\titlerule[0.4pt]]
\titlespacing*{\chapter}{0pt}{0pt}{16pt}

\titleformat{\section}
    {\normalfont\sffamily\Large\bfseries\color{ocpnavy}}
    {\thesection}{1em}{}
\titleformat{\subsection}
    {\normalfont\sffamily\large\bfseries\color{ocpblue}}
    {\thesubsection}{1em}{}
\titleformat{\subsubsection}
    {\normalfont\sffamily\normalsize\bfseries\color{darkslate}}
    {\thesubsubsection}{1em}{}

% --- Hyperliens & Métadonnées PDF ---
\usepackage[
    colorlinks=true,
    linkcolor=ocpnavy,
    citecolor=ocpgreen,
    urlcolor=ocpblue,
    bookmarks=true,
    bookmarksnumbered=true,
    pdfstartview=FitH
]{hyperref}

% --- Symboles & Outils Utiles ---
\usepackage{amsmath, amssymb}
\usepackage{enumitem}
\setlist{topsep=2pt, itemsep=2pt, parsep=1pt}

% Boîtes d'encadrés personnalisées
\usepackage[framemethod=TikZ]{mdframed}
\mdfdefinestyle{infobox}{
    innertopmargin=6pt,
    innerbottommargin=6pt,
    innerleftmargin=10pt,
    innerrightmargin=10pt,
    backgroundcolor=ocpnavy!4,
    linecolor=ocpnavy,
    linewidth=1pt,
    roundcorner=3pt
}
